% Teil II: durchgerechneter Weg, Teile a--c.

\subsection*{Gegeben und gesucht}
Nach $n=100$ Perioden ist die Amplitude auf die Hälfte abgesunken
$x(100T)=\tfrac12 x(0)$; Schwingungsfrequenz $f=\qty{50}{Hz}$.
Gesucht: (a) $\Lambda$, (b) $\beta$, (c) Zeit bis $x=0.7\,x_0$.

\subsection*{Ansatz}
Amplitudenverhältnis pro Periode: $x(t+T)=e^{-\Lambda}x(t)$; über $n$ Perioden
also Faktor $e^{-n\Lambda}$. Aus ``Hälfte nach 100'' folgt $100\Lambda=\ln 2$.
Mit $T=1/f$ erhält man $\beta=\Lambda/T$; der 70\,\%-Zeitpunkt folgt aus
$e^{-\beta t}=0.7$.

\subsection*{Rechnung}
\textbf{a) Logarithmisches Dekrement.}
\begin{equation}
  x(100T)=\tfrac12 x(0)
  \;\Rightarrow\;
  e^{-100\Lambda}=\tfrac12
  \;\Rightarrow\;
  \boxed{\;\Lambda=\frac{\ln 2}{100}\approx6.93\times10^{-3}\;}.
\end{equation}

\textbf{b) Dämpfungskonstante.} Mit $T=1/f=\qty{0.02}{s}$:
\begin{equation}
  \beta=\frac{\Lambda}{T}=\Lambda\, f
  =\frac{6.93\times10^{-3}}{0.02}
  \approx \boxed{\;\qty{0.347}{s^{-1}}\;}.
\end{equation}

\textbf{c) Zeit für 70\,\%.} $x(t)=0.7\,x_0\Rightarrow e^{-\beta t}=0.7$:
\begin{equation}
  t=-\frac{\ln 0.7}{\beta}
  =\frac{0.3567}{0.3466}
  \approx\boxed{\;\qty{1.03}{s}\;}.
\end{equation}

\begin{ergebnis}
$\Lambda=\ln2/100\approx6.93\times10^{-3}$, $\beta=\qty{0.347}{s^{-1}}$,
$t_{0.7}\approx\qty{1.03}{s}$.
\end{ergebnis}

\begin{kontrolle}
\textbf{Einheiten:} $\beta=\Lambda/T$ ist \unit{\per\second}; $t=\ln 0.7/\beta$
in \unit{\second}. \checkmark

\textbf{Querprobe:} Nach $t\approx1\,\mathrm s$ ($\approx50$ Perioden) ist die
Amplitude $e^{-0.3466}\approx0.71$ --- knapp über 70\,\%. Nach 100 Perioden
$e^{-0.693}\approx0.50$ --- passt zur Aufgabenangabe. \checkmark
\end{kontrolle}

\begin{merksatz}
Das Logarithmische Dekrement ist die Dämpfung pro Periode ($\Lambda=\beta T$);
die Amplitude fällt über $n$ Perioden um den Faktor $e^{-n\Lambda}$, die
Energie um $e^{-2n\Lambda}$.
\end{merksatz}
